\documentclass[10pt,a4paper]{amsart}
\usepackage[margin=19mm]{geometry}
\usepackage{amsmath,amssymb,mathtools}
\usepackage[scr=rsfs,cal=euler]{mathalfa}
\usepackage{xfrac}
\usepackage[unicode,hidelinks,bookmarks=false]{hyperref}
\newcommand{\ie}{i.e.\ }
\newcommand{\cf}{cf.\ }
\newcommand{\resp}{respectively\ }
\newcommand{\Subsection}{\subsection}
\providecommand{\nobreakdash}{\nobreak\hbox{-}}
\setlength{\emergencystretch}{2em}
\multlinegap0pt
\usepackage{xcolor}
\usepackage{enumerate}
\usepackage[shortlabels]{enumitem}
\usepackage{amssymb}
\usepackage{tikz}
\newtheorem{conjecture}{Conjecture}
\newtheorem{proposition}{Proposition}
\def\Pb{\mathbb{P}}
\def\bR{\mathbb{R}}
\title{On the Probability a Weighted Bernoulli Sum Exceeds Its Mean}
\author{Aleksa Milojevi\'c, Benny Sudakov}
\date{Source: arXiv:2606.30287v1 (CC BY 4.0)}
\begin{document}
\maketitle

\noindent\emph{Source and licence.} On the Probability a Weighted Bernoulli Sum Exceeds Its Mean. Version arXiv:2606.30287v1 (CC BY 4.0), distributed by arXiv under the Creative Commons Attribution 4.0 International licence, http://creativecommons.org/licenses/by/4.0/, as recorded in the arXiv licence metadata for this exact version and shown on the abstract page https://arxiv.org/abs/2606.30287v1. Full licence notice: \texttt{licenses/arxiv-2606.30287-CC-BY-4.0.txt}.\par\smallskip

\noindent\textit{Editorial scope.} Counted unit: the article's proposition prop:counterexample (source0.tex lines 210--212). The article's complete original proof of it is delivered with it (source0.tex lines 213--220, one \texttt{proof} environment of 8 lines). No mathematics is written, repaired or re-derived: the statement and the proof are the article's own lines, in the article's own notation, and no retained line is substituted or re-flowed.  Retained uncounted as the source-local context the proof needs: lines 132--135.  Nothing was omitted from any retained span, so the package carries no omission record.  No retained line contains a citation, so the wrapper carries no bibliography and no bibliography entry.  No retained line contains a figure environment, an image-inclusion command or drawing code, so no image is delivered and the package declares no asset dependency.  Editorial formatting only, in addition: an AMS \texttt{amsart} preamble that restates the article's own declarations and macros used by the retained lines, namely the environments \texttt{conjecture}, \texttt{proposition} on separate counters, together with the standard packages those lines need.  The retained environments print in this document as Conjecture 1, Proposition 1 in order of appearance, whereas the article prints the same items with the numbering of its own full document.  The complete native source of the article as distributed in the arXiv e-print for this exact version is bundled unchanged as \texttt{sources/204008/source0.tex}.
\begin{conjecture}\label{conj:prob}
Let $0\leq p\leq 1/3$ be a fixed real number. Let $w\in \bR_{>0}^m$ be such that $\sum_{i=1}^m w_i=1$, and let $v\in \{0, 1\}^m$ be a random vector with i.i.d. Bernoulli$(p)$ coordinates. Then
\[\Pb[\langle v, w\rangle\geq p]\geq p.\]
\end{conjecture}

\begin{proposition}\label{prop:counterexample}
Conjecture~\ref{conj:prob} does not hold for any $p\in (1/3, 1)$, with the exception of $p=1/2$. 
\end{proposition}

\begin{proof}
Let us first consider the case $1/3<p<1/2$. In this case, choosing $w=(1/3, 1/3, 1/3)$, we see that $\langle w, v\rangle\geq p$ if and only if $v$ has ones on at least $2$ coordinates. Thus,
\[\Pb[\langle w, v\rangle\geq p]=3p^2(1-p)+p^3=p(3p-2p^2).\]
It is not hard to verify that this quantity is smaller than $p$ for all $p\in (1/3, 1/2)$, since 
\[p-p(3p-2p^2)=p(1-3p+2p^2)=p(1-2p)(1-p)>0.\]
If $1/2<p<1$, we have a similar example, with $w=(1/2, 1/2)$. Again, we see that $\langle w, v\rangle\geq p$ if and only if $v$ has ones on at least $2$ coordinates. Thus,
\[\Pb[\langle w, v\rangle\geq p]=p^2<p.\qedhere\]
\end{proof}

\end{document}
